\documentclass[10pt,a4paper]{amsart}
\usepackage[margin=19mm]{geometry}
\usepackage{amsmath,amssymb,mathtools}
\usepackage[scr=rsfs,cal=euler]{mathalfa}
\usepackage{xfrac}
\usepackage[unicode,hidelinks,bookmarks=false]{hyperref}
\newcommand{\ie}{i.e.\ }
\newcommand{\cf}{cf.\ }
\newcommand{\resp}{respectively\ }
\newcommand{\Subsection}{\subsection}
\providecommand{\nobreakdash}{\nobreak\hbox{-}}
\setlength{\emergencystretch}{2em}
\multlinegap0pt
\usepackage{bm}
\usepackage{amssymb}
\usepackage{algorithm}\usepackage{algpseudocode}
\newtheorem{proposition}{Proposition}
\title{\LARGE \bf Moral Hazard in LTI Dynamics: A Hypothesis Testing Approach}
\author{Jaewon Jeong, Pan-Yang Su, S. Shankar Sastry,, Anil Aswani}
\date{Source: arXiv:2605.00158v1 (CC BY 4.0)}
\begin{document}
\maketitle

\noindent\emph{Source and licence.} \LARGE \bf Moral Hazard in LTI Dynamics: A Hypothesis Testing Approach. Version arXiv:2605.00158v1 (CC BY 4.0), distributed by arXiv under the Creative Commons Attribution 4.0 International licence, http://creativecommons.org/licenses/by/4.0/, as recorded in the arXiv licence metadata for this exact version and shown on the abstract page https://arxiv.org/abs/2605.00158v1. Full licence notice: \texttt{licenses/arxiv-2605.00158-CC-BY-4.0.txt}.\par\smallskip

\noindent\textit{Editorial scope.} Counted unit: the article's proposition prop:ytdist (source0.tex lines 315--318). The article's complete original proof of it is delivered with it (source0.tex lines 320--328, one \texttt{proof} environment of 9 lines). No mathematics is written, repaired or re-derived: the statement and the proof are the article's own lines, in the article's own notation, and no retained line is substituted or re-flowed.  Retained uncounted as the source-local context the proof needs: lines 246--255.  Nothing was omitted from any retained span, so the package carries no omission record.  No retained line contains a citation, so the wrapper carries no bibliography and no bibliography entry.  No retained line contains a figure environment, an image-inclusion command or drawing code, so no image is delivered and the package declares no asset dependency.  Editorial formatting only, in addition: an AMS \texttt{amsart} preamble that restates the article's own declarations and macros used by the retained lines, namely the environments \texttt{proposition} on separate counters, together with the standard packages those lines need.  The retained environments print in this document as Proposition 1 in order of appearance, whereas the article prints the same items with the numbering of its own full document.  The complete native source of the article as distributed in the arXiv e-print for this exact version is bundled unchanged as \texttt{sources/199542/source0.tex}.
\begin{equation}
\label{eqn:uvw}
    U_i = \begin{bmatrix}
C & 0 & \dots & 0 \\
C(A+BK_i) & C & \dots & 0 \\
C(A+BK_i)^2 & C(A+BK_i) & \dots & 0 \\
\vdots & \vdots & \ddots & \vdots \\
C(A+BK_i)^{T-1} & C(A+BK_i)^{T-2} & \dots & C
\end{bmatrix}\qquad  V_i = \begin{bmatrix} C(A+BK_i)^{\hphantom{1}} \\ C(A+BK_i)^2 \\ C(A+BK_i)^3 \\ \vdots \\ C(A+BK_i)^T \end{bmatrix}
\end{equation}

\begin{proposition}
\label{prop:ytdist}
    Under the state-feedback controller $K_i$ for $i = 0$ or $i = 1$, we have $\mathbf{Y}_T \sim \mathcal{N}(M_i, S_i)$ for $M_i = U_i\boldsymbol{\mu}_W + V_i\mu_0$ and $S_i = U_i\mathbf{\Sigma}_WU_i^\top + V_i\Sigma_0V_i^\top + \mathbf{\Sigma}_E$, with $U_i\in\mathbb{R}^{(T\cdot p)\times(T\cdot n)}$ and $V_i \in \mathbb{R}^{(T\cdot p)\times n}$ given in (\ref{eqn:uvw}), and where $\boldsymbol{\mu}_W \in \mathbb{R}^{T\cdot n}$ is defined by stacking $\mu_w$ a total of $T$ times, $\mathbf{\Sigma}_W = \textrm{blkdiag}(\Sigma_w, \ldots, \Sigma_w) \in \mathbb{R}^{(T\cdot n)\times(T\cdot n)}$ where $\textrm{blkdiag}(\cdot)$ specifies a block diagonal matrix, and $\mathbf{\Sigma}_E = \textrm{blkdiag}(\Sigma_e, \ldots,\Sigma_e) \in \mathbb{R}^{(T\cdot p)\times(T\cdot p)}$.
\end{proposition}

\begin{proof}
    Let $\mathbf{W}_T \in \mathbb{R}^{T\cdot n}$ be defined by stacking the random variables $w_0,\ldots,w_{T-1}$, and let $\mathbf{E}_T \in \mathbb{R}^{T\cdot p}$ be defined by stacking the random variables $\epsilon_0,\ldots,\epsilon_{T-1}$. Under the assumptions in our model, it is well known that $x_{t} = (A+BK_i)^tx_0 + \sum_{\tau = 0}^{t-1}(A+BK_i)^{t-\tau-1}w_\tau$, implying
    \begin{equation}
        \mathbf{Y}_T = U_i\mathbf{W}_T + V_ix_0 + \mathbf{E}_T.\label{eqn:lincom}
    \end{equation} 
    Since $x_0, \mathbf{W}_T, \mathbf{E}_T$ are jointly Gaussian and independent, they are all together jointly Gaussian. Hence, any linear combination such as (\ref{eqn:lincom}) is also jointly Gaussian. This means we completely know the distribution of $\mathbf{Y}_T$ if we calculate its mean and covariance. By linearity of expectation, we have $\mathbb{E}(\mathbf{Y}_T) = U_i\mathbb{E}(\mathbf{W}_T) + V_i\mathbb{E}(x_0) + \mathbb{E}(\mathbf{E}_T) = 
        U_i\boldsymbol{\mu}_W + V_i\mu_0$. If we denote the covariance using the notation $\textrm{Cov}(\cdot)$, then $\textrm{Cov}(\mathbf{Y}_T) = \textrm{Cov}(U_i\mathbf{W}_T) + \textrm{Cov}(V_ix_0) + \textrm{Cov}(\mathbf{E}_T) = 
        U_i\mathbf{\Sigma}_WU_i^\top + V_i\Sigma_0V_i^\top + \mathbf{\Sigma}_E$, where we used independence of $x_0,\mathbf{W}_T,\mathbf{E}_T$ in the first equality. The result follows.
\end{proof}

\end{document}
